\section{Conclusion}\label{sec:min-conclusion}

C-TreePO pairs a small algebra with a design-based audit. The algebra
states three local laws on a recursive summarization tree: each leaf
preserves the task evidence its span carries, re-summarizing a stored
summary does not drift, and split-then-merge preserves the same cross-span
interactions as summarize-together. The audit samples realized leaves,
stored summaries, and merge nodes with logged propensities, runs an oracle
check on each, and returns IPW estimates of violation rates and oracle
distortion. Together they make a long-document measurement claim
auditable: a fixed calibration set sets the judge-versus-oracle residual
once, and any combination of full-document or sampled-node labels
tightens the remaining estimation envelope along a $C/\sqrt{n+1}$ curve.

On the Manifesto Project benchmark, a single open-weight C-tree
pipeline equals or outperforms a frontier-ensemble baseline on six
policy dimensions, with leaf-size invariance as the headline
structural result: the macro correlation against expert means stays
stable across a $16\times$ range of leaf sizes. The same invariance pattern reappears in
controlled simulations (Markov changepoints, HyperLogLog parity,
mergeable-sketch sims), confirming that the algebra is tracking a real
structural phenomenon rather than a stylized property of one DGP.

The framework's scope is locally-decomposable oracles. It does not claim
to certify generic long-context LLM use, and the audit reports honest
violations when the oracle does not localize. Inside that scope the
practical contribution is to make the preservation claim observable.
Root accuracy, local-law distortion, judge calibration, and sampling
uncertainty are separate objects rather than one hidden engineering step,
and the same audit certifies the gold-to-strong-to-small distillation
chain when published results land. The Semantic Forests companion takes
these ideas to orchestration and deployment at system scale.
